\section{Level 5: Model-Based Planning and Control}
\label{sec:level5}

\subsection{Capability Definition and Adjacent Boundary}

Level~5 uses a learned or mechanistic world model operationally to rank, select, optimize, evaluate, or control actions. The model must participate in the decision loop: it may generate candidate rollouts, provide a simulated environment, support model-predictive control or tree search, evaluate policies, or train an agent through imagined experience. A direct model-free policy and a paper that mentions planning only as future work do not meet this contract \citep{SuttonBarto2018Reinforcement}.

The adjacent boundary is functional rather than causal. Level~3 simulates a future after an action, and Level~4 estimates alternative outcomes under an identification contract. Level~5 adds operational action selection. It does not guarantee Level~4. A surgical controller can be trained in a simulator without estimating patient counterfactual outcomes. An offline sepsis planner can optimize a learned environment whose action-response estimates remain confounded. Conversely, a careful counterfactual estimator can remain at Level~4 if it never selects or ranks actions. Calling Level~5 the highest capability does not make it the strongest evidence category.

This distinction matters because optimization amplifies model weaknesses. A predictor is evaluated on states observed in the data; a planner actively searches for actions and trajectories that may exploit errors or leave empirical support. Small transition biases can compound across the horizon, and reward misspecification can direct the policy toward clinically undesirable states. Planning is thus a demanding use of a world model, not a shortcut to clinical actionability.

\subsection{Clinical Evidence Landscape}

\subsubsection{Clinical Policy Learning and Treatment Planning}

Treatment planning and model-based clinical reinforcement learning are the most direct Level~5 settings. Early sepsis work learned patient transitions for policy optimization or built explicit simulation environments, while neighboring approaches optimized treatment policies directly from observational data \citep{Raghu2018ModelBasedReinforcement,Parbhoo2018Improvingcounterfactualreasoning,Komorowski2018AIClinician,Kiani2019SepsisWorldModel}. DeepSurv and related recommenders illustrate a distinct boundary: ranking treatments from estimated outcomes is decision-relevant, but it is not automatically model-based planning over a transition system \citep{Katzman2018DeepsurvPersonalizedTreatment}. We include such studies because they reveal what a Level~5 claim must add beyond risk-adjusted treatment ranking.

Benchmarks and safety-oriented offline methods make this distinction testable. DTR-Bench supplies an in silico environment for dynamic treatment regimes, whereas inverse-constrained reinforcement learning addresses safety-critical objectives \citep{Luo2024DtrBenchSilico,Fang2024OfflineInverseConstrained}. Matching-based off-policy evaluation focuses on whether a proposed ventilation policy can be compared credibly from retrospective data \citep{Lee2025MatchingBasedOff}. These components do not all constitute medical world models, but they define the evaluation infrastructure and alternatives against which model-based planning should be judged.

Recent systems make the dynamics-to-decision connection more explicit. CLARITY, medDreamer, a general Medical World Model, and patient-dynamics agents use latent trajectories, imagined experience, or model interaction to support clinical choices \citep{Ding2025CLARITYMedicalWorld,Yang2025MedicalWorldModel2,Xu2025medDreamer,Wu2026AgentifyingPatientDynamics}. Physiology-informed intervention simulation, knowledge-guided population policy modeling, and treatment-response optimization extend this pattern across ECG, opioid-overdose policy, and clinical decision support \citep{Chen2026ECGWMPhysiology,Ma2026Policy4oodKnowledgeGuided,Qin2026TreatmentResponseOptimized}. In these systems, the model is inspectable inside the decision loop. Their claims remain bounded by retrospective support, reward design, and whether simulated action effects are reliable where the optimizer searches.

\subsubsection{Metabolic, Rehabilitation, and Public-Health Control}

Metabolic care provides one of the clearest translational pathways because glucose dynamics, sensing, and interventions can be represented at short time scales. Digital-twin-enabled precision nutrition, reinforcement-learning glycemic control, and adaptive artificial-pancreas systems span retrospective development, proof-of-concept clinical evaluation, and closed-loop control \citep{Shamanna2021Type2Diabetes,Wang2023Optimizedglycemiccontrol,Kovatchev2025InteractiveArtificialPancreas}. Their evidence is not interchangeable: a clinical outcome associated with a digital-twin program does not isolate the twin's causal contribution, whereas closed-loop control places stronger demands on calibration, fail-safes, and continual state updating.

Model-based control has also been proposed for rehabilitation and systemic inflammation \citep{Larie2022PreparingnextCOVID,Ye2025TowardsAIbased}. These settings broaden the action space from drug dosing to multimodal therapy and contextual adaptation. They also make reward validity harder. Functional recovery, inflammatory regulation, adherence, and delayed adverse outcomes cannot be reduced safely to a single short-horizon surrogate without sensitivity analysis.

\subsubsection{Mechanistic Digital Twins and Model-Predictive Control}

Mechanistic and hybrid twins can connect explicit physiology to optimization. The UVA/Padova diabetes simulator is a longstanding example of a virtual-patient environment used to develop and test control strategies \citep{Man2014UvaPadovaType}. Patient-specific optimization has also been studied for trauma surgery, radiotherapy under uncertainty, ventricular-tachycardia ablation, Alzheimer disease treatment, sepsis control, and acute-kidney-injury interventions \citep{Aubert2021DevelopmentDigitalTwins,Chaudhuri2023Predictivedigitaltwin,Chrispin2026DigitalTwinGuided,Xu2026NeuralOperatorBased,Pickard2026EhrMpcInference,Zhang2026WorldModelEnhanced}. Across these domains, the decisive evidence is not that the twin can be optimized, but that the recommended action remains stable under plausible parameter, structural, and observation uncertainty.

Level~5 is therefore not simply a synonym for reinforcement learning. Model-predictive control may use a mechanistic simulator, a learned generative patient model, or a hybrid transition and repeatedly replan as measurements arrive. Search may compare a finite set of treatment schedules without learning a policy. Both satisfy the capability contract when the model operationally ranks actions; neither establishes clinical validity without decision-level evaluation.

\subsubsection{Surgical, Acquisition, and Physiological Control}

Surgical and acquisition systems provide another route because actions are instrumented and short-horizon outcomes can be observed at high temporal resolution. World models have been used for general surgical grasping, visuomotor grasping, video-based policy learning, and open simulation for augmented dexterity \citep{Lin2024WorldModelsGeneral,unknown2025VisuomotorGraspingWorld,unknown2025SurgWorldCosmosH,Yu2024OrbitSurgicalOpen}. Related work applies goal-conditioned control to ultrasound navigation and reinforcement-learning path planning for robotic imaging \citep{Amadou2024GoalConditionedReinforcement,Bi2026AutonomousPathPlanning}. These tasks make planning technically natural, but image-space success does not demonstrate safe tissue interaction or reliable acquisition of diagnostically sufficient views.

Endovascular navigation illustrates the progression from a learned environment to a safety-constrained autonomy claim. Work on mechanical-thrombectomy navigation, online evaluation of surgical robot policies, and world-model-based endovascular intervention combines simulation, hierarchy, or policy assessment \citep{Robertshaw2025WorldModelAi,Zbinden2026CosmosSurgDVRK,Robertshaw2026SafeAutonomousRobotic,Robertshaw2026AiAutonomousNavigation}. The relevant validation extends beyond trajectory completion to collision, force, anatomical variability, recovery from error, and operator takeover. A model can support tool-motion planning or policy training in a simulator \citep{Yue2025EchoWorldLearningMotion,Rapuri2026SAWTowardSurgical}; clinical use still requires safety variables that may be absent from video realism metrics.

\subsubsection{Virtual Clinical Environments and Diagnostic Agents}

Virtual clinical environments can also train or evaluate diagnostic agents. Evolving diagnostic agents in a simulated clinic is relevant to Level~5 when the environment changes in response to tests, questions, or diagnostic actions and the agent uses those transitions to select what to do next \citep{unknown2025EvolvingDiagnosticAgents}. The key boundary is whether the environment represents a patient process or merely scores a completed answer. Diagnostic success alone does not demonstrate a world model; sequential action-dependent state change does.

Mechanistic and digital-twin systems can be optimized through model-predictive control, search, or parameterized treatment planning. Their appeal is interpretability: actions and physical or physiological constraints can be represented explicitly. Their limitation is model mismatch. A patient-specific twin built from sparse data may produce precise recommendations that are highly sensitive to uncertain parameters. Repeated state updating and uncertainty-aware planning are therefore central to the credibility of this pathway.

Level~5 also includes candidate-policy evaluation and simulator-mediated agent training when the world model is used operationally rather than as background motivation. These uses can be valuable before clinical deployment. A simulator may support rehearsal, data augmentation, stress testing, or controller pretraining even when it is not valid enough to recommend real care. The intended use determines the required fidelity: training environment, offline evaluation, and bedside decision support impose different error tolerances.

\subsection{Representative Model Designs}

Learned latent environments coupled to actor--critic algorithms or search are common model-based designs. The world model predicts latent transitions and outcomes; the planner evaluates candidate actions through imagined trajectories. Model-predictive control repeatedly replans as new observations arrive, limiting the horizon over which errors accumulate. Tree search can explore discrete interventions, while hierarchical policies can divide long clinical decisions into slower strategic and faster operational actions.

Agent systems may query a patient-dynamics model as a tool. This makes the world-model contribution visible if the agent uses simulated futures to compare actions and if ablations isolate the effect of the dynamics model. An LLM that generates a recommendation from text without an explicit transition remains outside Level~5, even when its prompt invokes a world model. Similarly, a model used only to shape reward or initialize a policy does not establish model-based planning unless it participates in action evaluation.

Mechanistic and hybrid planners use ODEs, finite-element systems, pharmacokinetic models, or rule-based environments, sometimes with learned residuals or parameter inference. They can encode hard constraints and known action pathways, which is attractive in physiology and surgery. Their decision validity depends on parameter uncertainty, structural adequacy, and how the planner behaves when observations conflict with the model. Hybrid systems should state which component supplies the state, transition, action response, reward, and uncertainty.

The planner and world model should be separable enough to inspect. Without this separation, a strong policy result cannot show whether the model learned useful dynamics, the policy exploited simulator artifacts, or the reward encoded the answer. Ablations comparing model-based and model-free variants, short- and long-horizon planning, and alternative transition models can reveal where performance originates.

\subsection{Evidence Quality and Failure Modes}

Evaluation should connect model error to decision error. In addition to predictive accuracy, studies should test compounding error across rollout depth, robustness to transition perturbations, uncertainty under unsupported actions, and sensitivity to reward or objective design. Comparisons should include clinician imitation, myopic rules, model-free baselines, and policies evaluated under the same assumptions. A planner should abstain or revert to a conservative policy when the state or action lies outside support.

Offline policy evaluation is development evidence, not prospective proof. Importance weighting, learned evaluators, or the same simulator used for training can favor the proposed policy when support is weak or model bias is shared. Semi-synthetic and mechanistic environments can provide known counterfactuals, but their validity is limited by the environment specification. External cohorts with different practice patterns, prospective silent deployment, workflow studies, and ultimately trials are needed as the intended use approaches clinical decision support.

Reward misspecification is a distinctive failure mode. Short-term physiological stabilization may conflict with mortality or recovery; image response may not align with survival or quality of life; video realism may not align with tissue safety. Optimization can exploit these mismatches more aggressively than prediction. Papers should therefore define the clinical utility represented by the reward, report alternative objectives, and examine whether the selected policy changes when surrogate assumptions are varied.

The characteristic overclaim is to treat model participation in action selection as evidence of causal or clinical validity. A Level~5 system can be computationally correct and clinically unreliable. The capability label should be paired with separate statements about intervention semantics, causal support, outcome relevance, and validation setting. This allows planning work to be recognized without equating a simulated improvement with patient benefit.

\subsection{Level Takeaway and Next Threshold}

Level~5 completes the computational decision loop, but it does not complete the clinical evidence loop. Its promise is greatest when the action space is explicit, the world model is inspectable, uncertainty constrains planning, and evaluation connects errors in simulation to errors in decisions. Its risk is greatest when optimization searches outside the support of retrospective data or exploits a surrogate outcome.

\levelsummary
{Learned, mechanistic, and hybrid world models are already used for treatment planning, policy evaluation, acquisition guidance, surgical simulation, and agent training.}
{A model participating in action selection does not establish causal or clinical validity. Planning can magnify transition bias, unsupported extrapolation, and reward misspecification.}
{The field must move beyond capability labels toward decision-level evidence: support-aware planning, calibrated abstention, external and prospective validation, and evaluation against patient-relevant outcomes.}
